% Lösungen Zone (zusammenbau v0.1)
\einheitenkopf[][Kennst du schon]{Das kennst du schon}

Zuordnung: 1 → die Datengrundlage aller Einheiten · 2 → die Verteilungen der Einheiten 1 und 2 entstehen aus mehrstufigen Experimenten · 3 → die Binomialformen der Einheiten 3 und 4 · 4 → gewichtete Summen und prozentuale Abweichungen der Einheit 1 · 5 → die Rückwärtsaufgaben der Einheit 2 · 6 → die Standardabweichung der Einheit 3 als Wurzel der Varianz · 7 → die Verteilungen der Einheiten 1 und 2 entstehen aus mehrstufigen Experimenten · 8 → die Verteilungen der Einheiten 1 und 2 entstehen aus mehrstufigen Experimenten

\erg{1}{a) $P(X \ge 2) = 0{,}9$ \quad b) $P(X = 1) = \frac{108}{360} = 0{,}3$}
\erg{2}{a) $P = \frac{1}{4}$ \quad b) $P = \frac{4}{6} \cdot \frac{3}{5} = \frac{2}{5}$}
\erg{3}{a) ja \quad b) $n = 40$, $p = 0{,}3$}
\erg{4}{a) $= 4$ \quad b) $= 2{,}7$}
\erg{5}{a) $x = 4$ \quad b) $n > 17{,}5$, also $n = 18$}
\erg{6}{a) $= 8$ \quad b) $= \sqrt{21} \approx 4{,}58$}
\erg{7}{a) Fehler: Lea zählt nur die Reihenfolge rot–weiß, weiß–rot fehlt; richtig $P = 2 \cdot \frac{2}{5} \cdot \frac{3}{5} = \frac{12}{25}$}
\erg{8}{a) $P = 2 \cdot \frac{1}{3} \cdot \frac{2}{3} = \frac{4}{9}$}
